% ── Appendix G: Prompt Templates ──
\section{Prompt Templates}
\label{app:prompts}

All prompt templates used in the study are documented below. Each template follows the same structure with variations only in the bias manipulation.

\subsection{Base Prompt Template}

The base (control) scoring prompt:
\begin{quote}\small
\noindent\texttt{Evaluate the following response to the instruction on a scale of 1--5,\\
where 1 is worst and 5 is best.\\
\#\#\# Instruction: [item]\\
\#\#\# Response: [response]\\
\#\#\# Score:}
\end{quote}

\subsection{Rubric Order Probe Variants}

\paragraph{Normal (control):}
\begin{quote}\small
\noindent\texttt{...on a scale of 1--5, where 1 is worst and 5 is best.}
\end{quote}

\paragraph{Reversed:}
\begin{quote}\small
\noindent\texttt{...on a scale of 1--5, where 1 is best and 5 is worst.}
\end{quote}

\paragraph{Random mapping:}
\begin{quote}\small
\noindent\texttt{...on a scale of 1--5, where 1 corresponds to 'average',\\
2 to 'excellent', 3 to 'worst', 4 to 'best', 5 to 'poor'.}
\end{quote}

\subsection{Score ID Probe Variants}

\paragraph{Numeric (control):}
\begin{quote}\small
\noindent\texttt{...Score (1--5):}
\end{quote}

\paragraph{Letter grades:}
\begin{quote}\small
\noindent\texttt{...Score (A--E, where A is best and E is worst):}
\end{quote}

\paragraph{Descriptive labels:}
\begin{quote}\small
\noindent\texttt{...Score (Poor, Below Average, Average, Good, Excellent):}
\end{quote}

\subsection{Reference Answer Probe Variants}

\paragraph{No exemplar (control):}
The model is prompted with only the instruction and response, as shown in the base template above.

\paragraph{Good exemplar:}
\begin{quote}\small
\noindent\texttt{Example of an excellent response: [good\_example]\\
Now evaluate the following response...}
\end{quote}

\paragraph{Poor exemplar:}
\begin{quote}\small
\noindent\texttt{Example of a poor response: [poor\_example]\\
Now evaluate the following response...}
\end{quote}

\subsection{Inference Configuration}

\begin{itemize}
    \item \textbf{Temperature:} 0.0 (greedy decoding) for all local models
    \item \textbf{Max tokens:} 5 (single numeric score)
    \item \textbf{Stop sequences:} None (model generates naturally)
    \item \textbf{Repeats:} 3 per condition (deterministic at temperature 0)
    \item \textbf{Timeout (OpenRouter):} 15 seconds per API call
\end{itemize}

\subsection{Score Parsing}

Scores are parsed from model output using the following logic:

\begin{enumerate}
    \item Extract the first integer 1--5 from the output (numeric condition)
    \item Map letter grades A=5, B=4, C=3, D=2, E=1 (letter condition)
    \item Map descriptive labels: Excellent=5, Good=4, Average=3, Below Average=2, Poor=1 (descriptive condition)
    \item If no valid score is found, flag as parse failure (excluded from analysis)
\end{enumerate}

The descriptive parser was less reliable for base models, which sometimes output non-standard text. This affected 11.1\% of descriptive condition judgments. Numeric and letter variants had $<$ 1\% parse failure rate.
